% GENERATED FILE. Edit canonical lessons, then run npm run generate:books.
% canonical-source-hash: fnv1a64:41edf3398bac5e80
% canonical-lessons: MR-C410-kinarpatti, MR-C410-khadi, MR-C410-dipgrh, MR-C410-dvipkalpa, MR-C410-caupati

\chapter{Along the Coast}
\label{ch:mr-along-the-coast}
\hlchaptermodality{410}

\begin{chapteropening}
\textbf{By the end of this chapter:} \emph{I can talk about a coastline, a creek, a lighthouse, a peninsula and a beachfront.}

Complete the last lesson of chapter 410: I can talk about a coastline, a creek, a lighthouse, a peninsula and a beachfront.
\end{chapteropening}

\section[\texorpdfstring{\mr{किनारपट्टी} (kinārpaṭṭī)}{किनारपट्टी (kinārpaṭṭī)}]{\mr{किनारपट्टी} (kinārpaṭṭī) — a coastline}

\label{lesson:MR-C410-kinarpatti}



\begin{quote}
Before the new one: say the Marathi for to grow rough, then the Marathi for to slosh about.
\end{quote}

\subsection*{You'll want to know: \mr{किनारपट्टी}}
\textbf{\mr{किनारपट्टी}} — \emph{kinārpaṭṭī} — \textquotedblleft{}a coastline\textquotedblright{}.

\begin{quote}
\mr{इथली} \mr{किनारपट्टी} \mr{खूप} \mr{सुंदर} \mr{आहे}.
\end{quote}

\begin{quote}
\emph{ithlī kinārpaṭṭī khūp sundar āhe.}
\end{quote}

\begin{quote}
The coastline here is very beautiful.
\end{quote}

Feminine.

Built on \mr{किनारा}, a shore, and \mr{पट्टी}, a strip: a long strip of shore.

\begin{cousinweb}
\textbf{\mr{किनारपट्टी}} (\emph{kinārpaṭṭī}) is a Marathi compound of \textbf{\mr{किनारा}}, \textquotedblleft{}a shore\textquotedblright{}, a loan usually traced to Persian \emph{kinār}, and \textbf{\mr{पट्टी}}, \textquotedblleft{}a strip\textquotedblright{}.
\end{cousinweb}

\begin{grammarlens}[title={What you've built}]
One more word for getting around.
\end{grammarlens}

\subsection*{Your turn}
\begin{itemize}
\raggedright
  \item \emph{Say it:} \emph{kinārpaṭṭī}
  \item \emph{Say it:} \emph{kinārpaṭṭī}, once more
  \item \emph{Say it:} \emph{ithlī kinārpaṭṭī khūp sundar āhe.} — The coastline here is very beautiful.
  \item \emph{Recall:} say the Marathi for to grow rough, then the Marathi for to slosh about, then say \emph{kinārpaṭṭī} again
\end{itemize}

\begin{tcolorbox}[breakable,colback=teal!4,colframe=teal!35!black,title={Before you move on}]
What does \mr{किनारपट्टी} mean? (\textquotedblleft{}A coastline\textquotedblright{}.) Say it once more.
\end{tcolorbox}
\section[\texorpdfstring{\mr{खाडी} (khāḍī)}{खाडी (khāḍī)}]{\mr{खाडी} (khāḍī) — a creek; an inlet}

\label{lesson:MR-C410-khadi}



\begin{quote}
Before the new one: say the Marathi for to slosh about, then the Marathi for a coastline.
\end{quote}

\subsection*{You'll want to know: \mr{खाडी}}
\textbf{\mr{खाडी}} — \emph{khāḍī} — \textquotedblleft{}a creek; an inlet\textquotedblright{}.

\begin{quote}
\mr{खाडीत} \mr{मच्छीमारांच्या} \mr{होड्या} \mr{उभ्या} \mr{आहेत}.
\end{quote}

\begin{quote}
\emph{khāḍīt macchīmārāñcyā hoḍyā ubhyā āhet.}
\end{quote}

\begin{quote}
Fishermen’s boats are standing in the creek.
\end{quote}

Feminine.

Where the sea runs inland between the hills of the coast.

\begin{cousinweb}
\textbf{\mr{खाडी}} (\emph{khāḍī}) is inherited: CDIAL lists Marathi \emph{khāḍī}, \textquotedblleft{}creek, inlet\textquotedblright{}, under \emph{*khaḍḍa}, \textquotedblleft{}hole, pit\textquotedblright{} (CDIAL 3790).
\end{cousinweb}

\begin{grammarlens}[title={What you've built}]
One more word for getting around.
\end{grammarlens}

\subsection*{Your turn}
\begin{itemize}
\raggedright
  \item \emph{Say it:} \emph{khāḍī}
  \item \emph{Say it:} \emph{khāḍī}, once more
  \item \emph{Say it:} \emph{khāḍīt macchīmārāñcyā hoḍyā ubhyā āhet.} — Fishermen’s boats are standing in the creek.
  \item \emph{Recall:} say the Marathi for to slosh about, then the Marathi for a coastline, then say \emph{khāḍī} again
\end{itemize}

\begin{tcolorbox}[breakable,colback=teal!4,colframe=teal!35!black,title={Before you move on}]
What does \mr{खाडी} mean? (\textquotedblleft{}A creek; an inlet\textquotedblright{}.) Say it once more.
\end{tcolorbox}
\section[\texorpdfstring{\mr{दीपगृह} (dīpgṛh)}{दीपगृह (dīpgṛh)}]{\mr{दीपगृह} (dīpgṛh) — a lighthouse}

\label{lesson:MR-C410-dipgrh}



\begin{quote}
Before the new one: say the Marathi for a coastline, then the Marathi for a creek.
\end{quote}

\subsection*{You'll want to know: \mr{दीपगृह}}
\textbf{\mr{दीपगृह}} — \emph{dīpgṛh} — \textquotedblleft{}a lighthouse\textquotedblright{}.

\begin{quote}
\mr{रात्री} \mr{दीपगृहाचा} \mr{दिवा} \mr{दूरवर} \mr{दिसतो}.
\end{quote}

\begin{quote}
\emph{rātrī dīpgṛhācā divā dūrvar disto.}
\end{quote}

\begin{quote}
At night the lighthouse lamp can be seen far away.
\end{quote}

Neuter.

\mr{दीप} is a lamp and \mr{गृह} a house: a lamp house.

\begin{cousinweb}
\textbf{\mr{दीपगृह}} (\emph{dīpgṛh}) is a modern word built from Sanskrit \emph{dīpa}, \textquotedblleft{}a lamp\textquotedblright{}, and \emph{gṛha}, \textquotedblleft{}a house\textquotedblright{} (Monier-Williams).
\end{cousinweb}

\begin{grammarlens}[title={What you've built}]
One more word for getting around.
\end{grammarlens}

\subsection*{Your turn}
\begin{itemize}
\raggedright
  \item \emph{Say it:} \emph{dīpgṛh}
  \item \emph{Say it:} \emph{dīpgṛh}, once more
  \item \emph{Say it:} \emph{rātrī dīpgṛhācā divā dūrvar disto.} — At night the lighthouse lamp can be seen far away.
  \item \emph{Recall:} say the Marathi for a coastline, then the Marathi for a creek, then say \emph{dīpgṛh} again
\end{itemize}

\begin{tcolorbox}[breakable,colback=teal!4,colframe=teal!35!black,title={Before you move on}]
What does \mr{दीपगृह} mean? (\textquotedblleft{}A lighthouse\textquotedblright{}.) Say it once more.
\end{tcolorbox}
\section[\texorpdfstring{\mr{द्वीपकल्प} (dvīpkalpa)}{द्वीपकल्प (dvīpkalpa)}]{\mr{द्वीपकल्प} (dvīpkalpa) — a peninsula}

\label{lesson:MR-C410-dvipkalpa}



\begin{quote}
Before the new one: say the Marathi for a creek, then the Marathi for a lighthouse.
\end{quote}

\subsection*{You'll want to know: \mr{द्वीपकल्प}}
\textbf{\mr{द्वीपकल्प}} — \emph{dvīpkalpa} — \textquotedblleft{}a peninsula\textquotedblright{}.

\begin{quote}
\mr{द्वीपकल्पाच्या} \mr{तीन} \mr{बाजूंना} \mr{समुद्र} \mr{असतो}.
\end{quote}

\begin{quote}
\emph{dvīpkalpācyā tīn bājūnnā samudra asto.}
\end{quote}

\begin{quote}
A peninsula has sea on three sides.
\end{quote}

Masculine.

\mr{द्वीप} is an island and -\mr{कल्प} means almost: almost an island.

\begin{cousinweb}
\textbf{\mr{द्वीपकल्प}} (\emph{dvīpkalpa}) is a modern word built from Sanskrit \emph{dvīpa}, \textquotedblleft{}an island\textquotedblright{}, and \emph{kalpa}, \textquotedblleft{}like, almost\textquotedblright{} (Monier-Williams): land that is almost an island.
\end{cousinweb}

\begin{grammarlens}[title={What you've built}]
One more word for getting around.
\end{grammarlens}

\subsection*{Your turn}
\begin{itemize}
\raggedright
  \item \emph{Say it:} \emph{dvīpkalpa}
  \item \emph{Say it:} \emph{dvīpkalpa}, once more
  \item \emph{Say it:} \emph{dvīpkalpācyā tīn bājūnnā samudra asto.} — A peninsula has sea on three sides.
  \item \emph{Recall:} say the Marathi for a creek, then the Marathi for a lighthouse, then say \emph{dvīpkalpa} again
\end{itemize}

\begin{tcolorbox}[breakable,colback=teal!4,colframe=teal!35!black,title={Before you move on}]
What does \mr{द्वीपकल्प} mean? (\textquotedblleft{}A peninsula\textquotedblright{}.) Say it once more.
\end{tcolorbox}
\section[\texorpdfstring{\mr{चौपाटी} (caupāṭī)}{चौपाटी (caupāṭī)}]{\mr{चौपाटी} (caupāṭī) — a beachfront; a seafront}

\label{lesson:MR-C410-caupati}



\begin{quote}
Before the new one: say the Marathi for a lighthouse, then the Marathi for a peninsula.
\end{quote}

\subsection*{You'll want to know: \mr{चौपाटी}}
\textbf{\mr{चौपाटी}} — \emph{caupāṭī} — \textquotedblleft{}a beachfront; a seafront\textquotedblright{}.

\begin{quote}
\mr{रविवारी} \mr{चौपाटीवर} \mr{खूप} \mr{गर्दी} \mr{असते}.
\end{quote}

\begin{quote}
\emph{ravivārī caupāṭīvar khūp gardī aste.}
\end{quote}

\begin{quote}
On Sunday the beachfront is very crowded.
\end{quote}

Feminine.

A town beach with food stalls and evening walkers.

\begin{cousinweb}
\textbf{\mr{चौपाटी}} (\emph{caupāṭī}) names an open town beach where people walk and eat in the evening; its origin is not settled.
\end{cousinweb}

\begin{grammarlens}[title={What you've built}]
One more word for getting around.
\end{grammarlens}

\subsection*{Your turn}
\begin{itemize}
\raggedright
  \item \emph{Say it:} \emph{caupāṭī}
  \item \emph{Say it:} \emph{caupāṭī}, once more
  \item \emph{Say it:} \emph{ravivārī caupāṭīvar khūp gardī aste.} — On Sunday the beachfront is very crowded.
  \item \emph{Recall:} say the Marathi for a lighthouse, then the Marathi for a peninsula, then say \emph{caupāṭī} again
\end{itemize}

\begin{tcolorbox}[breakable,colback=teal!4,colframe=teal!35!black,title={Before you move on}]
What does \mr{चौपाटी} mean? (\textquotedblleft{}A beachfront; a seafront\textquotedblright{}.) Say it once more.
\end{tcolorbox}
